%
% 11-expnentialfunktion.tex -- Die Exponentialfunktion
%
% (c) 2026 Prof Dr Andreas Müller
%
\subsection{Die Exponentialfunktion}
Wir betrachten die Differentialgleichung
\begin{equation}
y' = -ky
\label{buch:differentialgleichungen:potenzreihen:eqn:expdgl}
\end{equation}
und setzen den
Potenzreihenansatz~\eqref{buch:differentialgleichungen:potenzreihen:eqn:ansatz}
in die Differentialgleichung ein.
Die Ableitung der  
Potenzreihen~\eqref{buch:differentialgleichungen:potenzreihen:eqn:ansatz}
ist
\[
y'(x)
=
\sum_{k=1}^\infty ka_kx^{k-1}.
\]
Einsetzen in die Differentialgleichung
\eqref{buch:differentialgleichungen:potenzreihen:eqn:expdgl}
ergibt die Gleichung
\begin{align*}
\sum_{l=1}^\infty a_llx^{l-1}
&=
-k
\sum_{l=0}\infty a_lx^l.
\intertext{Um diese besser vergleichen zu können, ersetzen wir $l-1$ durch
$l$, oder $l$ durch $l+1$ auf der linken Seite und erhalten}
\sum_{l=0}^\infty a_{l+1} (l+1)x^{l}
&=
-k
\sum_{l=0}^\infty a_lx^l.
\end{align*}
Durch Vergleich der Koeffizienten der beiden Potenzreihen können wir
schliessen, dass
\begin{equation}
a_{l+1}
=
-k
\frac{1}{l+1} a_l
\label{buch:differentialgleichungen:potenzreihen:eqn:exprek}
\end{equation}
ist für alle $l\ge 0$.
\eqref{buch:differentialgleichungen:potenzreihen:eqn:exprek}
ist eine Rekursionsgleichung, die ermöglicht, die Koefizienten $a_k$
mit $k>0$ auf $a_0$ zurückzuführen.
Löst man die Gleichung iterativ auf, erhält man die Werte
\begin{equation}
\begin{aligned}
a_1
&=
-k\frac{1}{0+1}a_0 = -ka_0 
\\
a_2
&=
-k\frac{1}{1+1}a_1 = -\frac{k}{2}a_1 = \phantom{-}\frac{k^2}{2}a_0
\\
a_3
&=
-k\frac{1}{2+1}a_2 = -\frac{k}{3}a_2 = -\frac{k^3}{3!}a_0
\\
a_4
&=
-k\frac{1}{3+1}a_3 = -\frac{k}{4}a_3 = \phantom{-}\frac{k^4}{4!}a_0
\\
a_5
&=
-k\frac{1}{4+1}a_4 = -\frac{k}{5}a_4 = -\frac{k^5}{5!}a_0
\\
&\hspace*{1mm}\vdots
\\
a_l
&=
(-1)^l \frac{k^l}{l!}a_0.
\end{aligned}
\label{buch:differentialgleichungen:potenzreihen:eqn:rekursion}
\end{equation}
Der Koeffizient $a_0$ wird durch die Differentialgleichung oder
durch \eqref{buch:differentialgleichungen:potenzreihen:eqn:expdgl}
nicht festgelegt.

Aus \eqref{buch:differentialgleichungen:potenzreihen:eqn:rekursion}
ergibt sich als Lösung der Differentialgleichung als die Potenzreihe
\[
y(x)
=
\sum_{l=0}^\infty a_lx^l
=
a_0
\sum_{l=0}^\infty
(-1)^l
\frac{k^l}{l!}x^l
=
a_0
\sum_{l=0}^\infty
\frac{(-kx)^l}{l!}
=
a_0e^{-kx}.
\]
Tatsächlich ist die Ableitung dieser Funktion
\[
y'(x)
=
-ka_0e^{-kx}
=
-ky(x),
\]
die Exponentialfunktion löst also die Differentialgleichung.

Da 
\eqref{buch:differentialgleichungen:potenzreihen:eqn:expdgl}
eine Differentialgleichung erster Ordnung ist, erwartet man 
genau eine frei wählbare Integrationskonstante.
Der Koeffizient $a_0$ ist diese Konstante.
Sie kann aus einer Anfangsbedingung $y(0)=a_0$ bestimmt werden.
